\begin{mistake}{2026-06-29}{06}

\begin{problemcontent}
设 \(a > 0\)，计算定积分：
\[
\int_0^{2a} x \sqrt{2ax - x^2} \mathrm{d}x
\]
\end{problemcontent}

\begin{solutioncontent}
对根号内二次式进行配方：
\[
2ax - x^2 = -(x^2 - 2ax + a^2 - a^2) = a^2 - (x-a)^2
\]
原积分化为：
\[
I = \int_0^{2a} x \sqrt{a^2 - (x-a)^2} \mathrm{d}x
\]
作平移换元，令 \(t = x - a\)，则 \(x = t + a\)，\(\mathrm{d}x = \mathrm{d}t\)。
积分限由 \([0, 2a]\) 变为 \([-a, a]\)：
\[
I = \int_{-a}^{a} (t + a) \sqrt{a^2 - t^2} \mathrm{d}t
\]
展开分为两项积分：
\[
I = \int_{-a}^{a} t \sqrt{a^2 - t^2} \mathrm{d}t + a \int_{-a}^{a} \sqrt{a^2 - t^2} \mathrm{d}t
\]
第一项被积函数 \(t \sqrt{a^2 - t^2}\) 为奇函数，在对称区间 \([-a, a]\) 上积分为 \(0\)。
第二项积分 \(\int_{-a}^{a} \sqrt{a^2 - t^2} \mathrm{d}t\) 的几何意义为半径是 \(a\) 的半圆面积，等于 \(\frac{1}{2}\pi a^2\)。
故：
\[
I = 0 + a \cdot \left( \frac{1}{2} \pi a^2 \right) = \frac{\pi}{2} a^3
\]
\end{solutioncontent}

\mistakeknowledge{
\begin{itemize}
  \item \textbf{核心公式}：对称区间奇偶性积分性质；半圆面积公式 \(\int_{-a}^{a} \sqrt{a^2 - t^2} \mathrm{d}t = \frac{1}{2}\pi a^2\)。
  \item \textbf{易错点}：带参数 \(a\) 配方时常数项易算错；不作平移换元直接算计算量极大且易错。
  \item \textbf{解题技巧}：遇到 \(\sqrt{2ax - x^2}\) 或区间为 \([0, 2a]\) 时，优先配方后作平移换元 \(t = x - a\)，构造 \([-a, a]\) 利用奇偶性及几何意义化简。
\end{itemize}}

\end{mistake}
